% Synthetic LaTeX article for tests/test_extract_pdf.py (compiled with XeLaTeX through
% h0lon.render.latex.compile_latex; the PDF next to it is committed for runners without TeX).
% Page 1: prose only; page 2: display and inline formulas (CMMI/CMSY/CMEX); page 3: prose
% with a single italic letter; a running header and page numbers on every page.
\documentclass[12pt]{article}
\usepackage[a4paper,margin=2cm,headheight=15pt]{geometry}
\usepackage{fontspec}
\IfFontExistsTF{Cambria}{\setmainfont{Cambria}}{\IfFontExistsTF{DejaVu Serif}{\setmainfont{DejaVu Serif}}{}}
\usepackage{amsmath,amssymb}
\usepackage{fancyhdr}
\pagestyle{fancy}
\fancyhf{}
\lhead{Численные методы. Тест-2026}
\rhead{\thepage}
\renewcommand{\headrulewidth}{0pt}
\begin{document}

\section*{Введение}

Этот короткий текст придуман для проверки извлечения и не взят ни из какого курса. Речь в
нём пойдёт о линейных системах: сколько операций стоит их решение, насколько точен
полученный ответ и как он меняется, если исходные числа слегка искажены. Почти всё здесь
держится на разложении матрицы в произведение более простых множителей.

Разложение достаточно построить один раз, а потом применять его к разным правым частям.
Ниже мы сравним прямой и итерационный подходы, заметим, что маленький остаток ещё не
гарантирует верного решения, и поговорим о том, зачем к задаче добавляют штрафное слагаемое.
В конце кратко опишем, как проверять ответы на заранее подготовленных примерах.

\newpage
\section*{Нормы и пределы}

Для $1 \le p < \infty$ определим векторную норму и норму Фробениуса матрицы $A = (a_{ij})$:
\[
  \|x\|_p = \Bigl(\sum_{i=1}^{n} |x_i|^p\Bigr)^{1/p}, \qquad
  \|A\|_F = \sqrt{\sum_{i,j} a_{ij}^2}.
\]
Пусть $A \in \mathbb{R}^{n \times n}$ невырождена и $f$ непрерывна на $[0, 1]$. Тогда
\[
  \int_0^1 f(t)\,dt = \lim_{n \to \infty} \frac{1}{n} \sum_{k=1}^{n} f\!\left(\frac{k}{n}\right),
  \qquad \kappa(A) = \|A\|\,\|A^{-1}\| \ge 1.
\]
Число обусловленности $\kappa_p(A)$ показывает, во сколько раз может вырасти относительная
ошибка решения по сравнению с ошибкой данных.

\newpage
\section*{Заключение}

Мы рассмотрели основные разложения и обсудили, когда решение системы устойчиво. Для
решения $x$ важно не только малое значение невязки, но и обусловленность задачи. В
следующей работе займёмся итерационными методами и сравним их скорость на больших
разреженных матрицах из практических задач.

\end{document}
